\scanpage{152}
supposées de degré impair, i.e. échangeant parité des degrés) pour des coalgèbres, comme correspondant aux dérivations resp. antidérivations de la $O_k$-Algèbre $V(M)$ (\unsure{évt|cat.} graduée mod 2). On a donc que le crochet de deux dérivations est une dérivation, l'anti-crochet de deux antidérivations est une antidérivation. On notera que si $\unsure{d}$ (revenant à un groupe de degrés quelc.) un endom. $d$ de $M$ est de degré $\gamma$, alors l'endom correspondant de $V(M)$ est de degré $-\gamma$ --- sauf si on change les degrés \inserted{de signe} pour l'objet
\diagram{\[
\begin{array}{ll}
V(M_i)&=V(M)^i=V(M)_{-i}\\
\text{\textquotedbl}&\\
V(M^{-i})&
\end{array}
\]}
Mais pour les histoires de dérivations \inserted{et antidérivations}\deleted{il est} seul un degré mod 2 importe, donc il n'y a pas alors à distinguer entre degré $i$ et $-i$ d'un endom. linéaire.

On traite de façon analogue les ``copuissances divisées'', \deleted{\unsure{un ``coidéal'' d'une coalgèbre}}\inserted{modules sur, sous-coalgèbre d'une} coalgèbre --- Pour le faire avec la généralité qui convient (pas seulement dans le cas commu., \deleted{\unsure{\&}} mais aussi le cas anticommu.) il faut d'abord disposer des \unsure{notions} dans le cas pas ``co'', mais ordinaire.
\margin{151}
